% Figure environments. All three are the UCI-HAR sweep, to match Section IV.
% Regenerate with `python analyze.py results_har/`, which writes the PNGs into
% results_har/; copy them beside main.tex (graphicspath is ./ only, on purpose
% -- results/ holds identically named CIFAR-10 figures).
% Inputted before Section IV in main.tex: floats only move forward.

\begin{figure}[t]
\centering
\includegraphics[width=\columnwidth]{fig1_convergence.png}
\caption{Central test accuracy against communication round on UCI-HAR over
100 rounds, means of three seeds with $\pm 1$ SD bands. The marker is the
energy-optimal stopping round, drawn only where accuracy ends more than half
a point below its peak: here $\varepsilon = 0.5$ alone, peaking at $0.4712$ at
round 17 and ending at $0.4553$. The remaining budgets are still flat or
rising at round 100, so the peak ordering CIFAR-10 shows across budgets does
not reproduce here; what does is that a tight enough budget makes further
rounds cost energy and accuracy together.}
\label{fig:convergence}
\end{figure}

% fig2 (energy to target) dropped: Table~\ref{tab:target} carries the same
% numbers with the unreached cells legible, which a bar chart cannot show.

\begin{figure}[t]
\centering
\includegraphics[width=\columnwidth]{fig3_energy_per_round.png}
\caption{Energy per round on UCI-HAR by privacy budget, net of idle, with the
non-private baseline in grey and $\pm 1$ SD over three seeds. Private training
costs $21.0$~J per round against $14.5$~J, an overhead of $+45\%$, and varies
by $8.2\%$ of its mean across the five private conditions --- inside the
$9.46\%$ measurement noise floor. This separates the two mechanisms: the step
is an implementation overhead independent of $\varepsilon$, so every
budget-dependent cost acts through the round count.}
\label{fig:perround}
\end{figure}

\begin{figure}[t]
\centering
\includegraphics[width=\columnwidth]{fig4_carbon.png}
\caption{Emissions for 1000 UCI-HAR training runs of the identical measured
workload, across six grid carbon intensities. The vertical distance between a
private curve and the non-private one is the carbon attributable to privacy:
$105$--$133$~gCO$_2$eq on Algeria's gas-dominated grid against
$5.6$--$7.1$~g in Norway, a factor of $18.8$ for the same computation and the
same guarantee.}
\label{fig:carbon}
\end{figure}
